\documentclass[12pt]{article}
\usepackage[margin=0.8in]{geometry}
\usepackage{amsmath}
\usepackage{xcolor}
\usepackage{tcolorbox}
\usepackage{enumitem}

\title{\textbf{Certainty Equivalent: Simple Explanation}}
\author{}
\date{}

\begin{document}
\maketitle

\section*{What is a Certainty Equivalent?}

\begin{tcolorbox}[colback=blue!5!white,colframe=blue!75!black]
\textbf{Simple Definition:}

The certainty equivalent (CE) is the \textbf{guaranteed amount of money} that makes you exactly as happy as taking a risky gamble.

\vspace{0.3cm}

In other words: "How much would I need for certain to make me indifferent to this lottery?"
\end{tcolorbox}

\subsection*{The Formula}

The CE is defined by:
\[
u(CE) = EU
\]

where:
\begin{itemize}
\item $u$ = your utility function
\item $CE$ = certainty equivalent (what we're solving for)
\item $EU$ = expected utility of the gamble
\end{itemize}

\textbf{In words:} The utility of getting CE for certain equals the expected utility of the gamble.

\section*{Simple Example}

You're offered a coin flip:
\begin{itemize}
\item Heads: win \$100
\item Tails: win \$0
\end{itemize}

If you're risk-neutral, your CE = \$50 (the expected value).

But if you're risk-averse, you'd say: "I'd be just as happy with \$40 guaranteed as I would be with this 50-50 gamble for \$100." Then your CE = \$40.

\vspace{0.3cm}

\textbf{Key insight:} The more risk-averse you are, the LOWER your certainty equivalent!

\newpage

\section*{CRRA Utility Example}

\subsection*{The Setup}

Let's say you use CRRA utility:
\[
u(x) = \frac{x^{1-\sigma}}{1-\sigma}
\]

where $\sigma$ is your risk aversion parameter (higher $\sigma$ = more risk averse).

\subsection*{The Gamble}

50-50 chance between:
\begin{itemize}
\item \$0 with probability $\frac{1}{2}$
\item \$200 with probability $\frac{1}{2}$
\end{itemize}

\subsection*{Finding the Certainty Equivalent}

\textbf{Step 1: Calculate Expected Utility}

\[
EU = \frac{1}{2} \cdot u(0) + \frac{1}{2} \cdot u(200)
\]

With CRRA utility:
\[
EU = \frac{1}{2} \cdot \frac{0^{1-\sigma}}{1-\sigma} + \frac{1}{2} \cdot \frac{200^{1-\sigma}}{1-\sigma}
\]

Since $0^{1-\sigma} = 0$:
\[
EU = \frac{1}{2} \cdot \frac{200^{1-\sigma}}{1-\sigma}
\]

\textbf{Step 2: Set $u(CE) = EU$}

\[
\frac{CE^{1-\sigma}}{1-\sigma} = \frac{1}{2} \cdot \frac{200^{1-\sigma}}{1-\sigma}
\]

\textbf{Step 3: Cancel out the $(1-\sigma)$ terms}

\[
CE^{1-\sigma} = \frac{1}{2} \cdot 200^{1-\sigma}
\]

\textbf{Step 4: Solve for CE}

Take both sides to the power of $\frac{1}{1-\sigma}$:
\[
CE = \left(\frac{1}{2}\right)^{\frac{1}{1-\sigma}} \cdot 200
\]

\begin{tcolorbox}[colback=yellow!20!white,colframe=orange!75!black]
\textbf{The Answer:}
\[
\boxed{CE = \left(\frac{1}{2}\right)^{\frac{1}{1-\sigma}} \cdot 200}
\]
\end{tcolorbox}

\newpage

\section*{What Does This Mean?}

Let's plug in different values of $\sigma$ (risk aversion):

\subsection*{Case 1: Risk Neutral ($\sigma = 0$)}

\[
CE = \left(\frac{1}{2}\right)^{\frac{1}{1-0}} \cdot 200 = \left(\frac{1}{2}\right)^1 \cdot 200 = \frac{1}{2} \cdot 200 = 100
\]

Makes sense! If you're risk-neutral, you only care about expected value:
\[
EV = \frac{1}{2}(0) + \frac{1}{2}(200) = 100
\]

\subsection*{Case 2: Slightly Risk Averse ($\sigma = 0.5$)}

\[
CE = \left(\frac{1}{2}\right)^{\frac{1}{1-0.5}} \cdot 200 = \left(\frac{1}{2}\right)^2 \cdot 200 = \frac{1}{4} \cdot 200 = 50
\]

You'd need \$50 guaranteed to be as happy as the gamble.

\subsection*{Case 3: More Risk Averse ($\sigma = 2$)}

\[
CE = \left(\frac{1}{2}\right)^{\frac{1}{1-2}} \cdot 200 = \left(\frac{1}{2}\right)^{-1} \cdot 200 = 2 \cdot 200 = 400
\]

Wait, that's weird! This says CE = \$400, but the most you can win is \$200!

\textbf{What's happening?} When $\sigma > 1$, the exponent $\frac{1}{1-\sigma}$ becomes negative, which flips the fraction. This is mathematically correct but the interpretation changes.

For $\sigma > 1$ with this particular gamble (one outcome is \$0), the utility function becomes very sensitive to getting zero, and the math doesn't give us the intuitive CE we expect.

\subsection*{Case 4: Log Utility ($\sigma = 1$)}

When $\sigma = 1$, we use $u(x) = \ln(x)$ instead:

\[
EU = \frac{1}{2} \ln(0) + \frac{1}{2} \ln(200)
\]

But $\ln(0) = -\infty$, so this gamble has infinite negative utility! You'd reject it at any price.

\begin{tcolorbox}[colback=red!5!white,colframe=red!75!black]
\textbf{Important Note:}

When one outcome is \$0 and you have log utility or high risk aversion, the gamble can have very low (or even infinitely negative) utility because you really, really don't want to end up with \$0.

That's why in practice, these problems usually give you gambles where all outcomes are positive amounts.
\end{tcolorbox}

\newpage

\section*{The Big Picture}

\subsection*{How Risk Aversion Affects CE}

\begin{center}
\begin{tabular}{|c|c|l|}
\hline
\textbf{Risk Aversion ($\sigma$)} & \textbf{CE} & \textbf{Interpretation} \\
\hline
0 (risk-neutral) & \$100 & CE = Expected Value \\
\hline
0.5 (slightly risk-averse) & \$50 & Need \$50 for sure \\
\hline
Higher $\sigma$ & Lower CE & Need even less for sure \\
\hline
\end{tabular}
\end{center}

\vspace{0.5cm}

\textbf{The pattern:} As you become MORE risk-averse (higher $\sigma$), your certainty equivalent DECREASES.

\vspace{0.3cm}

\textbf{Intuition:} "I really hate risk, so I'll accept a much smaller guaranteed amount instead of gambling."

\subsection*{Risk Premium}

The \textbf{risk premium} is how much you're willing to give up to avoid risk:

\[
\text{Risk Premium} = E[x] - CE
\]

For our example with $\sigma = 0.5$:
\[
RP = 100 - 50 = 50
\]

You'd pay up to \$50 to avoid this risk!

\section*{More Realistic Example}

Let's do one where all outcomes are positive:

\subsection*{Gamble: 50-50 between \$100 and \$200}

Expected value: $E[x] = \frac{1}{2}(100) + \frac{1}{2}(200) = 150$

\subsection*{Finding CE with CRRA utility}

\[
EU = \frac{1}{2} \cdot \frac{100^{1-\sigma}}{1-\sigma} + \frac{1}{2} \cdot \frac{200^{1-\sigma}}{1-\sigma}
\]

\[
\frac{CE^{1-\sigma}}{1-\sigma} = \frac{1}{2} \cdot \frac{100^{1-\sigma}}{1-\sigma} + \frac{1}{2} \cdot \frac{200^{1-\sigma}}{1-\sigma}
\]

Cancel $(1-\sigma)$:
\[
CE^{1-\sigma} = \frac{1}{2} \left(100^{1-\sigma} + 200^{1-\sigma}\right)
\]

Solve for CE:
\[
CE = \left[\frac{1}{2} \left(100^{1-\sigma} + 200^{1-\sigma}\right)\right]^{\frac{1}{1-\sigma}}
\]

\subsection*{Numerical Examples}

\textbf{Risk-neutral ($\sigma = 0$):}
\[
CE = \left[\frac{1}{2}(100 + 200)\right]^1 = \frac{300}{2} = 150
\]
CE = Expected Value, as expected!

\textbf{Risk-averse ($\sigma = 0.5$):}
\[
CE = \left[\frac{1}{2}(100^{0.5} + 200^{0.5})\right]^{2} = \left[\frac{1}{2}(10 + 14.14)\right]^2 \approx (12.07)^2 \approx 145.7
\]

So you'd accept \$145.70 guaranteed instead of the 50-50 gamble for \$100 or \$200.

Risk premium: $150 - 145.7 = \$4.30$

\newpage

\section*{Summary}

\begin{tcolorbox}[colback=green!5!white,colframe=green!75!black]
\textbf{Key Takeaways:}

\begin{enumerate}[leftmargin=*]
\item \textbf{Certainty Equivalent (CE):} The guaranteed amount that makes you as happy as taking the gamble
\[
u(CE) = EU
\]

\item \textbf{Finding CE:}
\begin{itemize}
\item Calculate expected utility $EU$
\item Solve $u(CE) = EU$ for $CE$
\end{itemize}

\item \textbf{Risk aversion and CE:}
\begin{itemize}
\item More risk-averse $\Rightarrow$ Lower CE
\item Risk-neutral $\Rightarrow$ CE = Expected Value
\end{itemize}

\item \textbf{Risk Premium:}
\[
RP = E[x] - CE
\]
How much you'd pay to avoid risk

\item \textbf{CRRA utility:} $u(x) = \frac{x^{1-\sigma}}{1-\sigma}$
\begin{itemize}
\item $\sigma$ = risk aversion parameter
\item Higher $\sigma$ = more risk-averse = lower CE
\end{itemize}

\item \textbf{Watch out:} When outcomes include \$0 and $\sigma \geq 1$, utility can be $-\infty$ (you really hate getting zero!)
\end{enumerate}
\end{tcolorbox}

\section*{Exam Tips}

\begin{itemize}[leftmargin=*]
\item \textbf{Always start with:} $u(CE) = EU$

\item \textbf{For CRRA:} The $(1-\sigma)$ denominators will cancel when you set $u(CE) = EU$

\item \textbf{Don't forget:} Take the $\frac{1}{1-\sigma}$ power at the end to solve for CE

\item \textbf{Sanity check:}
\begin{itemize}
\item CE should be between the minimum and maximum outcomes (usually)
\item CE $\leq$ Expected Value (if risk-averse)
\item CE = Expected Value (if risk-neutral)
\end{itemize}

\item \textbf{Special case:} If $\sigma = 1$, use log utility: $u(x) = \ln(x)$
\end{itemize}

\end{document}
